
\documentclass{article}
\usepackage{colortbl}
\usepackage{makecell}
\usepackage{multirow}
\usepackage{supertabular}

\begin{document}

\newcounter{utterance}

\centering \large Interaction Transcript for game `textmapworld', experiment `large\_cycle', episode 41 with reka{-}flash{-}3.1{-}t0.0.
\vspace{24pt}

{ \footnotesize  \setcounter{utterance}{1}
\setlength{\tabcolsep}{0pt}
\begin{supertabular}{c@{$\;$}|p{.15\linewidth}@{}p{.15\linewidth}p{.15\linewidth}p{.15\linewidth}p{.15\linewidth}p{.15\linewidth}}
    \# & $\;$A & \multicolumn{4}{c}{Game Master} & $\;\:$B\\
    \hline

    \theutterance \stepcounter{utterance}  
    & & \multicolumn{4}{p{0.6\linewidth}}{
        \cellcolor[rgb]{0.9,0.9,0.9}{
            \makecell[{{p{\linewidth}}}]{
                \texttt{\tiny{[P1$\langle$GM]}}
                \texttt{Please help me with the following task. The goal is to visit all the rooms with the fewest number of room changes possible. In each room, you need to decide the direction to go in. Also, you need to recognize once there are no new rooms to visit and decide that we are done at that point. Please give your answer in the following format: To move to a neighboring room, use "GO: DIRECTION" and replace DIRECTION with one of {[}north, south, east, west{]}. To stop the exploration, answer with "DONE" instead. Omit any other text.} \\
\texttt{Here is an example:} \\
\texttt{You are in the Kitchen. Currently available directions: south, west. What is your next instruction?} \\
\texttt{GO: west} \\
\texttt{You have made a step and entered a Lobby. Currently available directions: east, north. What is your next instruction?} \\
\texttt{GO: north} \\
\texttt{...} \\
\texttt{You have made a step and entered a Bedroom. Currently available directions: south. What is your next instruction?} \\
\texttt{DONE} \\
\texttt{Let us start. You are in the Recreation room. Currently available directions: south, west. What is your next instruction?} \\
            }
        }
    }
    & \\ \\

    \theutterance \stepcounter{utterance}  
    & \multicolumn{4}{p{0.6\linewidth}}{
        \cellcolor[rgb]{0.9,0.9,0.9}{
            \makecell[{{p{\linewidth}}}]{
                \texttt{\tiny{[P1$\rangle$GM]}}
                \texttt{<reasoning>} \\
\texttt{Let's tackle this problem. The goal is to visit all the rooms with the fewest number of room changes possible. So, I need to figure out the optimal path that covers all rooms in the least steps. But how do I know the layout of the rooms? The example starts in the Kitchen and moves west to the Lobby, then north to a Bedroom, etc. But in the actual problem, the user starts in the Recreation room with available directions south and west.} \\
\\ 
\texttt{Wait, but the user hasn't provided the entire map. So maybe this is a dynamic problem where each step reveals new rooms and their directions. The challenge is to navigate through the rooms, keeping track of visited ones to avoid cycles, and find the shortest path to cover all rooms. But since the problem is presented step by step, perhaps each time the user gives a direction, the next room is revealed, and I have to make decisions based on that.} \\
\\ 
\texttt{Wait, the example shows that each time the user gives a direction, the assistant responds with the next instruction. So the assistant is supposed to guide the user through the rooms, making decisions at each step. But the assistant needs to plan the path to visit all rooms with minimal room changes. Wait, but the problem says "the goal is to visit all the rooms with the fewest number of room changes possible." So room changes refer to the number of times you switch directions, i.e., the number of steps taken. Because each "GO: DIRECTION" is a step, and each step is a room change. So the total number of steps (room changes) should be minimized to visit all rooms.} \\
\\ 
\texttt{But how do we know when all rooms are visited? The user says "recognize once there are no new rooms to visit and decide that we are done at that point." So when the current room has no unvisited neighboring rooms, and all rooms have been visited, then we stop. Wait, but how do we track which rooms have been visited? Because each time we enter a room, we need to mark it as visited. But in the problem's interaction, the assistant doesn't have memory of previous rooms. Wait, but in the example, the assistant starts in the Kitchen, goes west to Lobby, then north to Bedroom, etc. So the assistant must be keeping track of visited rooms as the conversation progresses.} \\
\\ 
\texttt{Wait, but in the actual problem, the user is the one providing the current room and available directions.} \\
            }
        }
    }
    & & \\ \\

    \theutterance \stepcounter{utterance}  
    & & & \multicolumn{2}{p{0.3\linewidth}}{
        \cellcolor[rgb]{0.9,0.9,0.9}{
            \makecell[{{p{\linewidth}}}]{
                \texttt{\tiny{[GM$|$GM]}}
                \texttt{abort game} \\
            }
        }
    }
    & & \\ \\

\end{supertabular}
}

\end{document}
